\section{Introduction}
Citizens increasingly turn to general-purpose large language models (LLMs) for
questions of law: \emph{Am I entitled to a year-end bonus? What is the legal
working week? Can my landlord keep the deposit?} The models reply in fluent,
authoritative prose---and sometimes cite articles that do not exist, that exist
but have been repealed, or that exist and are in force yet say something
entirely different from what the answer claims. In open-domain conversation
such slips are tolerable; in law they are not. The domain has an
\emph{asymmetric} cost structure: a confident, plausible-sounding, fabricated
citation actively \emph{misleads} a citizen who cannot easily distinguish it
from a real one, and is therefore strictly worse than an honest admission of
ignorance. Recent work profiling \emph{legal} hallucinations reports that
general-purpose LLMs invent or misattribute legal authorities at high
rates~\cite{dahl2024legalfictions}, confirming that the problem is systemic
rather than anecdotal.

Hallucination in natural language generation is by now well
documented~\cite{ji2023hallucination, huang2025hallucination}, and
retrieval-augmented generation (RAG) is the standard
mitigation~\cite{lewis2020rag, gao2023ragsurvey}: the model is conditioned on
retrieved passages so that its output is, ideally, grounded in real text. Yet
grounding \emph{reduces} but does not \emph{eliminate} invented references. A
model handed the correct passage can still transpose a digit in an article
number, cite a neighbouring article, or attribute a real quotation to the wrong
code. Grounding improves the \emph{expected} faithfulness of a draft; it does
not \emph{guarantee} any individual citation. In law, the guarantee is what
matters.

\paragraph{Our stance: refusing is a feature.}
\emph{Le Rapporteur} adopts a deliberately conservative position suited to the
domain, organized around a single, non-negotiable invariant: \textbf{no
verifiable source, no answer}. The invariant is enforced not by prompt
engineering---advisory, and easily undone by the model's own sampling---but by
a \emph{post-generation verification stage} that is \emph{fail-closed}. The
model drafts freely; but before anything reaches the user, every citation in
the draft is checked, one by one, against an authoritative legal database. If a
single citation cannot be resolved to a real, in-force article, or if any error
occurs while checking (a network timeout, a malformed reference, an ambiguous
code), the system does not guess: it returns an explicit
refusal---\emph{``Je ne trouve pas de texte applicable.''} (\emph{I find no
applicable text.}). Doubt collapses to refusal, never to fabrication. This is
defense in depth: the trust guarantee is a property of the \emph{system} and
holds even when the model misbehaves.

\paragraph{A running example.}
Two questions recur throughout the paper. Asked \emph{``Do I have a right to a
Christmas bonus?''}, a fluent model confidently invents
\emph{article~L.\,4321-7 du Code du travail}---an authority that simply does
not exist. Asked for \emph{``the legal working week''}, a live model we tested
answered \emph{article~L.\,3132-2} (which exists, but governs weekly
\emph{rest}) instead of the correct \emph{L.\,3121-27}. The first is a
non-existent authority; the second is real-but-wrong. A system that merely
pattern-matches ``looks like a citation'' passes both; Le Rapporteur refuses
the first outright and, as we discuss (\S\ref{sec:eval},~\S\ref{sec:limits}),
exposes the second as the harder, semantic case.

\paragraph{Contributions.}
\begin{itemize}
  \item A \textbf{fail-closed pipeline} for legal QA---\emph{generate,
  extract, verify, source-or-refuse}---whose trust guarantee is a property of
  the \emph{system}, not of the model, stated as a precise release predicate
  (\S\ref{sec:overview}).
  \item A \textbf{database-centric citation verifier} over the
  Canutes/Légifrance corpus, and the query engineering---index-friendly
  variant expansion, accent-folded code disambiguation, state-aware selection
  of the in-force version---that makes per-citation checking feasible at
  interactive latency over a large, irregular schema without full-table
  regular-expression scans (\S\ref{sec:verify}).
  \item A \textbf{reusable trust service}: the verifier is exposed over HTTP
  (\texttt{/verify}) and as a Model Context Protocol (MCP) server, so any
  third-party agent can fact-check the citations in an arbitrary LLM output
  without adopting our model (\S\ref{sec:service}).
  \item A \textbf{hardware-sovereign serving layer}: a thin OpenAI-compatible
  indirection lets the identical pipeline run on non-NVIDIA accelerators; we
  demonstrate it live on a Qualcomm Cloud~AI~100 and target AMD via a portable
  runtime (\S\ref{sec:sovereign}).
  \item An \textbf{executable trust specification} in Gherkin, readable and
  auditable by a jurist, in which a single fabricated article turns the build
  red (\S\ref{sec:bdd}).
\end{itemize}

\paragraph{Scope and honesty.}
This system was built under hackathon time constraints. We are deliberately
candid about what is \emph{proven} (runs today, tested) versus
\emph{envisioned} (designed, on the roadmap), and we mark the boundary
explicitly rather than blur it---an editorial choice mirroring the system's own
fail-closed ethos. The remainder of the paper follows the pipeline:
\S\ref{sec:background} frames the corpus and a taxonomy of citation errors;
\S\ref{sec:overview} presents the pipeline and its release predicate;
\S\ref{sec:verify} details the database verifier that is the technical core;
\S\ref{sec:service} the reusable trust service; \S\ref{sec:sovereign} the
sovereign serving layer; \S\ref{sec:bdd} the executable specification; and
\S\S\ref{sec:eval}--\ref{sec:limits} our observations and the road from
existence checking to semantic verification.
